\chapter{Your First Scan}
\label{ch:quickstart}

\newthought{The fastest way to understand \Jarvis} is to run it. This chapter takes you from
an empty directory to a finished scan with structured output, using the quick-start cards
that ship inside every new project. You will not write any \textsc{yaml} yet---that comes in
Part~\ref{part:task-card}---you will just run, watch, and read the results.

\section{Create a project}
\label{sec:create-project}

\Jarvis\ works inside a \emph{standalone project}: a self-contained directory that holds your
task cards and collects all outputs. Create one with:

\begin{shellbox}
Jarvis project create MyScan
cd MyScan
\end{shellbox}

\noindent This scaffolds a ready-to-run project:

\begin{jfiletree}
\dirtree{%
.1 \dtfold{MyScan/}.
.2 \dtfold{bin/}.
.3 \dtcode{quickstart\_mcmc\_operas.yaml}.
.3 \dtcode{quickstart\_csv\_operas.yaml}.
.2 \dtfold{data/}.
.3 \dtcsv{points.csv}.
.2 \dtfold{deps/}.
.3 \dtcode{environment\_default.yaml}.
.2 \dtdoc{README.md}.
.2 \dtmark{.jarvis-project.json}\DTcomment{project marker}.
.2 \dtyaml{jarvis.project.yaml}\DTcomment{project marker}.
}
\end{jfiletree}

\begin{description}[style=nextline, leftmargin=1.2em]
  \item[\dtfold{bin/}] Your \textsc{yaml} task cards live here; two quick-starts are provided.
  \item[\dtfold{data/}] External data files (here, a small \file{points.csv} for the replay
        example).
  \item[\dtfold{deps/}] Dependency declarations, including a default environment file.
  \item[\dtmark{.jarvis-project.json}, \dtyaml{jarvis.project.yaml}] \emph{Marker files} that tell
        \Jarvis\ ``this directory is the project root.'' They are what the \marker{\&J/} path
        marker resolves to (Chapter~\ref{ch:core-concepts}).
\end{description}

\begin{jnote}
Runtime directories---\file{outputs/}, \file{logs/}, and \file{images/}---are \emph{not}
created up front. \Jarvis\ makes them on first use, so a fresh project stays tidy.
\end{jnote}

\section{The quick-start card at a glance}
\label{sec:quickstart-card}

Open \file{bin/quickstart\_mcmc\_operas.yaml}. It is a complete, minimal task card
(shown below). Do not worry about every key yet---just match it to the
three-step mental model from Chapter~\ref{ch:introduction}.

\begin{yamlwin}[bin/quickstart\_mcmc\_operas.yaml]
Scan:
  name: "quickstart_mcmc_operas"
  save_dir: "&J/outputs"
  sample_directory:
    limit: 200
    width: 6
    archive_samples: true

Sampling:
  Method: "MCMC"
  Variables:
    - name: x
      description: "Toy variable x"
      distribution: {type: Flat, parameters: {min: 0, max: 5.0}}
    - name: y
      description: "Toy variable y"
      distribution: {type: Flat, parameters: {min: 0, max: 5.0}}
  Bounds:
    num_chains: 8
    num_iters: 600
    proposal_scale: 0.06
  LogLikelihood:
    - {name: "LogL_Z", expression: "LogGauss(z, 100, 10)"}

Operas:
  make_paraller: 8
  Modules:
    - name: EggBox
      operator: "helper.eggbox2d"
      call_mode: call
      required_modules: []
      input:
        - {name: x, expression: "x"}
        - {name: y, expression: "y"}
      output:
        - {name: z, entry: z}
\end{yamlwin}

\noindent Reading it against the three steps:

\begin{enumerate}
  \item \textbf{Draw a point.} \yamlkey{Sampling} runs the \sampler{MCMC} sampler over two
        flat variables \yamlkey{x} and \yamlkey{y} on $[0, 5]$.
  \item \textbf{Run a workflow.} \yamlkey{Operas} evaluates the in-process operator
        \code{helper.eggbox2d} on \yamlkey{x} and \yamlkey{y} and produces the observable
        \yamlkey{z}.
  \item \textbf{Score it.} \yamlkey{LogLikelihood} computes \expr{LogGauss(z, 100, 10)}---a
        Gaussian centred on a pretend measurement $z = 100 \pm 10$.
\end{enumerate}

\noindent The \yamlkey{Scan} block names the run and points outputs at \marker{\&J/outputs}.
(The card also declares an \yamlkey{EnvReqs} block, omitted from the listing for brevity, that
checks the platform before starting.)

\section{Run it}
\label{sec:run-it}

From the project root:

\begin{shellbox}
Jarvis bin/quickstart_mcmc_operas.yaml
\end{shellbox}

\noindent You will see, in order:

\begin{enumerate}
  \item the \Jarvis\ banner and version---the logging system is up;
  \item \code{Validation successful}---your card passed the schema check;
  \item the sampler initialising (here, \sampler{MCMC});
  \item the worker factory building and the workflow flowchart being drawn;
  \item periodic progress lines as samples are submitted and completed;
  \item an end-of-run summary from the \textsc{hdf5} writer and the factory, followed by an
        automatic conversion of the results to \textsc{csv}.
\end{enumerate}

\begin{jtip}
The terminal shows only the high-level status of the sampler and the factory. You do not need
to read it closely---everything important is captured in the project's \file{logs/} directory
and the end-of-run summary.
\end{jtip}

\section{Replay tabulated points}
\label{sec:replay}

Sometimes you already have the parameter points and just want to push them through the
workflow. The second quick-start does exactly that, reading rows from
\file{data/points.csv} instead of sampling:

\begin{shellbox}
Jarvis bin/quickstart_csv_operas.yaml
\end{shellbox}

\noindent Its only difference from the quick-start card above is the \yamlkey{Sampling}
block, which selects the \sampler{CSV} method and names the file and columns to replay:

\begin{yamlwin}
Sampling:
  Method: "CSV"
  CSV:
    path: "&J/data/points.csv"
    uuid_column: "uuid"
    variables: ["x", "y"]
  LogLikelihood:
    - {name: "LogL_Z", expression: "LogGauss(z, 100, 10)"}
\end{yamlwin}

\section{Look at what you produced}
\label{sec:first-outputs}

After either run, your project has gained a predictable set of artifacts under the run name:

\begin{jfiletree}
\dirtree{%
.1 \dtfold{MyScan/}.
.2 \dtfold{outputs/quickstart\_mcmc\_operas/}.
.3 \dtfold{DATABASE/}.
.4 \dtdata{samples.0.hdf5}\DTcomment{structured samples}.
.4 \dtcsv{samples.0.csv}\DTcomment{CSV export}.
.4 \dtcode{samples.schema.json}\DTcomment{observable schema}.
.3 \dtfold{SAMPLE/}\DTcomment{per-sample artifacts}.
.3 \dtcode{run\_summary.\{json,csv,txt\}}\DTcomment{run summaries}.
.2 \dtfold{logs/quickstart\_mcmc\_operas/}\DTcomment{run logs}.
.2 \dtfold{images/quickstart\_mcmc\_operas/}\DTcomment{flowchart.json, plots}.
}
\end{jfiletree}

\noindent The \file{DATABASE/} folder is what you analyse: \textsc{hdf5} for tools, \textsc{csv}
for a quick look, and a schema describing the columns. The \file{run\_summary.*} files report
how the run went. All of this is covered in detail in Part~\ref{part:running}.

\section{Convert, plot, and monitor}
\label{sec:quickstart-utilities}

Three optional commands operate on a run; each is detailed in Chapter~\ref{ch:cli}:

\begin{description}[style=nextline, leftmargin=1.2em]
  \item[\cli{Jarvis bin/quickstart\_mcmc\_operas.yaml -{}-convert}] Snapshot the live
        \textsc{hdf5} into \textsc{csv} while the scan is still running.
  \item[\cli{Jarvis bin/quickstart\_mcmc\_operas.yaml -{}-plot}] Generate a \JPlot\
        configuration from the task card (then render it with \cmd{jplot}).
  \item[\cli{Jarvis bin/quickstart\_mcmc\_operas.yaml -{}-monitor}] Open a real-time resource
        monitor for a running scan.
\end{description}

\section{The success checklist}
\label{sec:success-check}

Your first scan worked if all of the following are true:

\begin{itemize}
  \item the command exited without an error;
  \item \file{outputs/<run>/DATABASE/} contains a \file{.hdf5} and a \file{.csv};
  \item \file{run\_summary.txt} reports completed samples and zero (or few) failures.
\end{itemize}

\section{Make it your own}
\label{sec:make-your-own}

You now have the full loop working end to end. The natural next steps:

\begin{itemize}
  \item Learn the vocabulary---projects, markers, tokens, and the output
        contract---in Chapter~\ref{ch:core-concepts}.
  \item Understand every block of the task card in Part~\ref{part:task-card}.
  \item Wrap your own external program as a calculator in Chapter~\ref{ch:calculators}, using
        the same black-box pattern the quick-starts follow.
\end{itemize}
